% =====================================================================
\chapter{Del perceptrón al perceptrón multicapa}
\label{tema:1}
% =====================================================================

Este tema presenta la unidad básica de una red neuronal, el perceptrón, su regla de aprendizaje y el teorema que garantiza su convergencia, y muestra por qué una sola unidad no basta: la función XOR no es linealmente separable. La solución, añadir una capa oculta, conduce al perceptrón multicapa. La exposición sigue la §6.1 de~\cite{gbc} y las §6.1--6.3 de~\cite{jm3}; la demostración del teorema de convergencia sigue la lección~3 de~\cite{cs4780}, basada en Novikoff~\cite{novikoff}, y el recuento de parámetros, las notas «Neural Networks Part~1» de~\cite{cs231n}.

\section{La unidad y el perceptrón}

\begin{definicion}[Unidad]\label{def:unidad}
Una \emph{unidad} (o \emph{neurona artificial}) con $n$ entradas está determinada por un vector de \emph{pesos} $w\in\R^{n}$, un \emph{sesgo} $b\in\R$ y una función $\phi\colon\R\to\R$, llamada \emph{función de activación}. Para una entrada $x\in\R^{n}$ calcula la \emph{preactivación} $z=w\cdot x+b$ y la \emph{activación} $a=\phi(z)$.
\end{definicion}

\begin{definicion}[Perceptrón]\label{def:perceptron}
El \emph{perceptrón} de Rosenblatt~\cite{rosenblatt} es la unidad con activación escalón:
\[
\hat y=\ind[\,w\cdot x+b>0\,]\in\{0,1\}.
\]
Su \emph{frontera de decisión} es el hiperplano $\{x\in\R^{n}:w\cdot x+b=0\}$.
\end{definicion}

El perceptrón es, por tanto, un clasificador lineal: asigna la clase~1 al semiespacio abierto $\{w\cdot x+b>0\}$ y la clase~0 a su complementario.

\begin{ejemplo}[La puerta AND]\label{ej:and}
Con $w=(1,1)$ y $b=-1{,}5$, la entrada $(1,1)$ da $1+1-1{,}5=0{,}5>0$ y produce~1; las entradas $(0,0)$, $(0,1)$ y $(1,0)$ dan $-1{,}5$, $-0{,}5$ y $-0{,}5$, y producen~0. El perceptrón calcula la conjunción lógica.
\end{ejemplo}

\begin{definicion}[Separabilidad lineal]
Un conjunto finito de pares $(x_i,y_i)$ con $x_i\in\R^{n}$ e $y_i\in\{0,1\}$ es \emph{linealmente separable} si existen $w\in\R^{n}$ y $b\in\R$ tales que $w\cdot x_i+b>0$ cuando $y_i=1$ y $w\cdot x_i+b<0$ cuando $y_i=0$.
\end{definicion}

\section{La regla de aprendizaje de Rosenblatt}

\begin{definicion}[Regla del perceptrón]\label{def:regla}
Sea $\eta>0$ una \emph{tasa de aprendizaje}. Se parte de $w=0$, $b=0$ y se recorren los ejemplos de forma cíclica. Para cada ejemplo $(x_i,y_i)$ se calcula $\hat y_i$ y se actualiza
\[
w\leftarrow w+\eta\,(y_i-\hat y_i)\,x_i,\qquad b\leftarrow b+\eta\,(y_i-\hat y_i).
\]
El algoritmo termina cuando una pasada completa no produce ningún error.
\end{definicion}

Como $y_i-\hat y_i\in\{-1,0,1\}$, los pesos solo cambian cuando el ejemplo está mal clasificado, y el cambio desplaza la preactivación en la dirección correcta: si $y_i=1$ y $\hat y_i=0$, la nueva preactivación de $x_i$ aumenta en $\eta(\norm{x_i}^{2}+1)>0$.

\begin{ejemplo}[Aprendizaje de AND]\label{ej:rosenblatt}
Con $\eta=1$, recorriendo las entradas en el orden $(0,0)$, $(0,1)$, $(1,0)$, $(1,1)$, la regla realiza~10 correcciones y la sexta pasada ya no contiene errores. El resultado final es $w=(2,1)$, $b=-2$, que clasifica bien las cuatro entradas: la preactivación vale $-2$, $-1$, $0$ y $1$, y solo la última es positiva.
\end{ejemplo}

Para enunciar el teorema de convergencia conviene absorber el sesgo y usar etiquetas simétricas. Sea $\tilde x_i=(x_i,1)\in\R^{n+1}$, $\theta=(w,b)$ y $s_i=2y_i-1\in\{-1,+1\}$. El ejemplo $i$ está bien clasificado si $s_i\,\theta\cdot\tilde x_i>0$, y la regla de la \cref{def:regla} se escribe $\theta\leftarrow\theta+\eta\,s_i\,\tilde x_i$ cada vez que se comete un error. En particular, toda actualización se produce en un ejemplo con $s_i\,\theta\cdot\tilde x_i\le 0$.

\begin{teorema}[Convergencia del perceptrón; Novikoff]\label{teo:novikoff}
Supóngase que existen $u\in\R^{n+1}$ con $\norm{u}=1$ y $\gamma>0$ tales que $s_i\,u\cdot\tilde x_i\ge\gamma$ para todo $i$, y sea $R=\max_i\norm{\tilde x_i}$. Entonces, cualquiera que sea el orden en que se presentan los ejemplos, el número total $M$ de actualizaciones de la regla del perceptrón cumple
\[
M\le\frac{R^{2}}{\gamma^{2}}.
\]
\end{teorema}

\begin{proof}
Sea $\theta_k$ el vector de parámetros tras $k$ actualizaciones, con $\theta_0=0$, y supóngase que la actualización $k+1$ se produce en el ejemplo $i$, de modo que $\theta_{k+1}=\theta_k+\eta s_i\tilde x_i$ y $s_i\,\theta_k\cdot\tilde x_i\le 0$.

\emph{Crecimiento de la proyección sobre $u$.} Se tiene $\theta_{k+1}\cdot u=\theta_k\cdot u+\eta\,s_i\,\tilde x_i\cdot u\ge\theta_k\cdot u+\eta\gamma$. Por inducción, $\theta_M\cdot u\ge M\eta\gamma$.

\emph{Crecimiento de la norma.} Se tiene
\[
\norm{\theta_{k+1}}^{2}=\norm{\theta_k}^{2}+2\eta\,s_i\,\theta_k\cdot\tilde x_i+\eta^{2}s_i^{2}\norm{\tilde x_i}^{2}\le\norm{\theta_k}^{2}+\eta^{2}R^{2},
\]
porque el término central es $\le 0$ y $s_i^{2}=1$. Por inducción, $\norm{\theta_M}^{2}\le M\eta^{2}R^{2}$.

\emph{Conclusión.} Por la desigualdad de Cauchy--Schwarz y $\norm{u}=1$, $M\eta\gamma\le\theta_M\cdot u\le\norm{\theta_M}\le\sqrt{M}\,\eta R$. Dividiendo por $\sqrt{M}\,\eta\gamma>0$ (si $M=0$ no hay nada que probar) se obtiene $\sqrt{M}\le R/\gamma$.
\end{proof}

\begin{observacion}
La cota no depende de $\eta$ ni del número de ejemplos, sino solo de la geometría de los datos: el cociente $R/\gamma$ mide lo «estrecho» que es el pasillo que separa las dos clases. Si los datos no son linealmente separables, no existe tal $u$ y la regla no termina nunca.
\end{observacion}

\section{El problema XOR}

La función XOR vale~1 en $(0,1)$ y $(1,0)$ y vale~0 en $(0,0)$ y $(1,1)$.

\begin{proposicion}\label{prop:xor}
El conjunto $\{((0,0),0),((0,1),1),((1,0),1),((1,1),0)\}$ no es linealmente separable. En consecuencia, ningún perceptrón calcula XOR.
\end{proposicion}

\begin{proof}
Supóngase que existen $w=(w_1,w_2)$ y $b$ que lo separan. De las entradas con salida~0 se obtiene $b<0$ y $w_1+w_2+b<0$; sumando, $w_1+w_2+2b<0$. De las entradas con salida~1 se obtiene $w_2+b>0$ y $w_1+b>0$; sumando, $w_1+w_2+2b>0$. Las dos desigualdades son incompatibles.
\end{proof}

Minsky y Papert~\cite{minsky} analizaron en 1969 esta y otras limitaciones de los perceptrones de una capa, y su libro contribuyó a frenar la investigación en redes neuronales durante más de una década~\cite[cap.~1]{gbc}.

\section{Una capa oculta resuelve XOR}

La salida consiste en componer dos capas. Sea $\ReLU(t)=\max(0,t)$ y considérese la red de la \cref{fig:xor}, tomada de~\cite[§6.1]{gbc}:
\begin{equation}\label{eq:xor}
h_1=\ReLU(x_1+x_2),\qquad h_2=\ReLU(x_1+x_2-1),\qquad \hat y=h_1-2h_2.
\end{equation}

\begin{figure}[h]
\centering
\begin{tikzpicture}[>=Stealth,
  nodo/.style={circle,draw,minimum size=8mm,inner sep=0pt},
  every label/.style={font=\small}]
  \node[nodo] (x1) at (0,1) {$x_1$};
  \node[nodo] (x2) at (0,-1) {$x_2$};
  \node[nodo] (h1) at (3,1) {$h_1$};
  \node[nodo] (h2) at (3,-1) {$h_2$};
  \node[nodo] (y) at (6,0) {$\hat y$};
  \draw[->] (x1) -- node[above,pos=0.3,font=\small]{$1$} (h1);
  \draw[->] (x1) -- node[above,pos=0.25,font=\small]{$1$} (h2);
  \draw[->] (x2) -- node[below,pos=0.25,font=\small]{$1$} (h1);
  \draw[->] (x2) -- node[below,pos=0.3,font=\small]{$1$} (h2);
  \draw[->] (h1) -- node[above,font=\small]{$1$} (y);
  \draw[->] (h2) -- node[below,font=\small]{$-2$} (y);
  \node[font=\small,above=1pt of h1] {sesgo $0$};
  \node[font=\small,below=1pt of h2] {sesgo $-1$};
  \node[font=\small,anchor=west] at (7.2,0.9) {\begin{tabular}{@{}cccc@{}}\toprule $x$ & $h_1$ & $h_2$ & $\hat y$\\\midrule $(0,0)$ & 0 & 0 & 0\\ $(0,1)$ & 1 & 0 & 1\\ $(1,0)$ & 1 & 0 & 1\\ $(1,1)$ & 2 & 1 & 0\\\bottomrule\end{tabular}};
\end{tikzpicture}
\caption{Red con una capa oculta de dos unidades ReLU que calcula XOR.}\label{fig:xor}
\end{figure}

En el espacio de la capa oculta, las cuatro entradas se transforman en $(0,0)$, $(1,0)$, $(1,0)$ y $(2,1)$: las dos entradas con salida~1 caen en el mismo punto, y la función lineal $h_1-2h_2$ separa las clases. La capa oculta no clasifica: cambia la \emph{representación} para que una función lineal pueda hacerlo. Aprender la representación a la vez que el clasificador es la idea que distingue el aprendizaje profundo de los modelos lineales del módulo~8.

\section{El perceptrón multicapa}

\begin{definicion}[Perceptrón multicapa]\label{def:mlp}
Sean $n_0,n_1,\dots,n_L$ enteros positivos. Un \emph{perceptrón multicapa} (MLP) de $L$ capas con anchuras $n_0\to n_1\to\dots\to n_L$ está formado por matrices $W_\ell\in\R^{n_{\ell-1}\times n_\ell}$, sesgos $b_\ell\in\R^{1\times n_\ell}$ y activaciones $\phi_\ell$, $\ell=1,\dots,L$. Para $x\in\R^{1\times n_0}$ se define $h^{(0)}=x$ y
\[
z^{(\ell)}=h^{(\ell-1)}W_\ell+b_\ell,\qquad h^{(\ell)}=\phi_\ell\bigl(z^{(\ell)}\bigr),\qquad \ell=1,\dots,L,
\]
con $\phi_\ell$ aplicada componente a componente. Las capas $1,\dots,L-1$ se llaman \emph{ocultas}. En un clasificador de $K$ clases, $n_L=K$, $\phi_L$ es la identidad y la salida $z^{(L)}$ es el vector de logits, que se convierte en $p=\softmax(z^{(L)})$.
\end{definicion}

\begin{figure}[h]
\centering
\begin{tikzpicture}[>=Stealth,
  caja/.style={draw,rounded corners=2pt,minimum height=8mm,inner sep=4pt,font=\small}]
  \node[caja] (x) {$x\in\R^{1\times n}$};
  \node[caja,right=6mm of x] (l1) {$xW_1+b_1$};
  \node[caja,right=6mm of l1] (a) {$\phi(\cdot)$};
  \node[caja,right=6mm of a] (l2) {$hW_2+b_2$};
  \node[caja,right=6mm of l2] (s) {$\softmax$};
  \node[caja,right=6mm of s] (p) {$p\in\R^{1\times K}$};
  \draw[->] (x)--(l1); \draw[->] (l1)--node[above,font=\scriptsize]{$z^{(1)}$}(a);
  \draw[->] (a)--node[above,font=\scriptsize]{$h$}(l2); \draw[->] (l2)--node[above,font=\scriptsize]{$z^{(2)}$}(s); \draw[->] (s)--(p);
\end{tikzpicture}
\caption{Perceptrón multicapa con una capa oculta de anchura $H$ y salida softmax.}\label{fig:mlp}
\end{figure}

Con un lote de $B$ ejemplos apilados como filas de una matriz $X\in\R^{B\times n_0}$, la misma fórmula $XW_\ell+b_\ell$ (con $b_\ell$ sumado a cada fila) calcula todas las preactivaciones de la capa con un único producto de matrices. Esa es la razón práctica de la convención de vectores fila.

\begin{proposicion}[Número de parámetros]\label{prop:parametros-mlp}
Un MLP de anchuras $n_0\to n_1\to\dots\to n_L$ tiene $\sum_{\ell=1}^{L}(n_{\ell-1}+1)\,n_\ell$ parámetros.
\end{proposicion}

\begin{proof}
La capa $\ell$ tiene $n_{\ell-1}n_\ell$ pesos en $W_\ell$ y $n_\ell$ sesgos en $b_\ell$, es decir, $(n_{\ell-1}+1)n_\ell$ parámetros; se suman las $L$ capas.
\end{proof}

\begin{ejemplo}\label{ej:mnist}
Un MLP $784\to 128\to 10$ para clasificar dígitos de $28\times 28$ píxeles (módulo~8) tiene $785\cdot 128+129\cdot 10=100\,480+1\,290=101\,770$ parámetros.
\end{ejemplo}

\begin{proposicion}[Colapso de las capas lineales]\label{prop:colapso}
Si todas las activaciones de un MLP son la identidad, existen $W\in\R^{n_0\times n_L}$ y $b\in\R^{1\times n_L}$ tales que $z^{(L)}=xW+b$ para toda entrada $x$. Además, $\operatorname{rg}W\le\min_{\ell}n_\ell$.
\end{proposicion}

\begin{proof}
Por inducción sobre $L$. Para $L=1$ es la definición. Si $h^{(L-1)}=xW'+b'$, entonces $z^{(L)}=(xW'+b')W_L+b_L=x(W'W_L)+(b'W_L+b_L)$ por la distributividad y la asociatividad del producto de matrices; basta tomar $W=W'W_L$ y $b=b'W_L+b_L$. Resulta $W=W_1W_2\cdots W_L$, y el rango de un producto no supera el de ninguno de sus factores, que es a lo sumo $\min(n_{\ell-1},n_\ell)$ para $W_\ell$ (módulo~3).
\end{proof}

La \cref{prop:colapso} explica por qué la no linealidad es imprescindible: sin ella, añadir capas no aumenta lo que la red puede representar, y una capa intermedia estrecha incluso lo reduce.

\begin{observacion}[Conexión con los LLM]
Cada bloque de un Transformer contiene un MLP de dos capas que se aplica a la representación de cada token por separado, con anchuras $d\to 4d\to d$ en el Transformer original~\cite{vaswani} y en GPT-2. Sin contar sesgos, tiene $d\cdot 4d+4d\cdot d=8d^{2}$ pesos, frente a los $4d^{2}$ de las cuatro proyecciones de la atención (módulo~13). Los modelos abiertos recientes, como LLaMA, sustituyen la primera capa por una unidad con puerta (\cref{tema:2}).
\end{observacion}

\begin{lecturas}
\item Goodfellow, Bengio y Courville~\cite{gbc}, §6.1 «Example: Learning XOR» y §6.6 «Historical Notes».
\item Jurafsky y Martin~\cite{jm3}, §6.1--6.3: unidades, el problema XOR y redes prealimentadas.
\item Cornell CS4780~\cite{cs4780}, lección~3 «The Perceptron»: algoritmo y demostración de la convergencia.
\item Stanford CS231n~\cite{cs231n}, «Neural Networks Part~1: Setting up the Architecture», apartados sobre la organización en capas y el tamaño de las redes.
\item Bishop y Bishop~\cite{bishopdl}, cap.~6 «Deep Neural Networks».
\end{lecturas}

\section*{Ejercicios del tema 1}
\addcontentsline{toc}{section}{Ejercicios del tema 1}

\begin{ejercicio}{1}{eje:1-lineal}
(a) ¿Cuántas capas «efectivas» tiene un MLP de diez capas sin funciones de activación? (b) Dé un ejemplo de red lineal de dos capas $2\to 1\to 2$ cuya función no pueda obtenerse con ninguna matriz $2\times 2$ de rango~2, y explique por qué añadir capas lineales puede reducir la capacidad expresiva pero nunca aumentarla.
\fuente{adaptado de Zhang \emph{et al.}~\cite{d2l}, §5.1, ejercicio~1.}
\end{ejercicio}

\begin{ejercicio}{1}{eje:1-xor}
Explique por qué ningún perceptrón calcula XOR y qué hace la capa oculta de la red~\eqref{eq:xor} para que una función lineal pueda separar las clases. Indique las coordenadas de las cuatro entradas en el espacio $(h_1,h_2)$.
\end{ejercicio}

\begin{ejercicio}{2}{eje:1-param}
Calcule el número de parámetros de un MLP $784\to 256\to 128\to 10$ e indique qué fracción corresponde a la primera capa.
\end{ejercicio}

\begin{ejercicio}{2}{eje:1-cs231n}
Una red tiene 3 entradas, una capa oculta de 4 unidades y 2 salidas; otra tiene 3 entradas, dos capas ocultas de 4 unidades y una salida. Calcule para cada una el número de pesos, de sesgos y de parámetros totales.
\fuente{adaptado de Stanford CS231n~\cite{cs231n}, «Neural Networks Part~1», apartado «Sizing neural networks».}
\end{ejercicio}

\begin{ejercicio}{2}{eje:1-or}
Aplique la regla del perceptrón (\cref{def:regla}) con $\eta=1$, partiendo de $w=(0,0)$ y $b=0$, a la función OR, recorriendo las entradas en el orden $(0,0)$, $(0,1)$, $(1,0)$, $(1,1)$. Indique cada corrección, el número total de correcciones y los parámetros finales.
\end{ejercicio}

\begin{ejercicio}{3}{eje:1-novikoff}
(a) Demuestre que, partiendo de $\theta=0$, la sucesión de predicciones de la regla del perceptrón no depende de $\eta>0$. (b) Para la función AND, con la notación del \cref{teo:novikoff}, tome $u$ proporcional a $(1,1,-1{,}5)$ y calcule $\gamma$, $R$ y la cota $R^{2}/\gamma^{2}$. Compárela con las 10 correcciones del \cref{ej:rosenblatt}.
\end{ejercicio}
